\documentclass{article}
\usepackage[utf8]{inputenc}
\usepackage[spanish]{babel}
\usepackage{amsmath}
\usepackage{enumerate}
\usepackage[margin=2cm]{geometry}

\title{Problemas seleccionados para asesoría}
\author{}
\date{}

\begin{document}
\maketitle


\section{Error Radial en el Elipse de Aterrizaje de un Rover}

Análisis de desempeño de navegación y los márgenes de error en las misiones a Marte (\textit{Modern Spacecraft Guidance, Navigation, and Control}, Capítulo sobre \textit{Entry, Descent, and Landing Systems}).


Cuando un rover es lanzado hacia la superficie marciana, el punto de impacto real rara vez coincide con el punto de referencia (target). Definimos la variable aleatoria continua $X$ como el error radial de aterrizaje, es decir, la distancia euclidiana (en kilómetros) desde el centro del elipse de aterrizaje designado hasta el punto de impacto real.

Debido a que los errores en los ejes ortogonales (Norte y Este) suelen modelarse como variables normales independientes con media cero y misma varianza, la distancia radial resultante $X$ sigue una Distribución de Rayleigh. La Función de Densidad de Probabilidad (FDP) del error radial es:
$$ f(x) = \begin{cases} \frac{x}{\sigma^2} e^{-x^2 / (2\sigma^2)} & \text{para } x \ge 0 \\ 0 & \text{en otro caso} \end{cases} $$
donde $\sigma$ es el parámetro de escala que representa la dispersión característica del sistema de navegación.

\begin{enumerate}[a)]
    \item Derive la Función de Distribución Acumulativa (FDA) del error radial $X$.
    \item En la literatura de Guiado, Navegación y Control (GNC), el CEP (Circular Error Probable) es la métrica estándar de precisión. Se define como el radio del círculo centrado en el target que contiene el 50\% de las probabilidades de impacto. Calcule el CEP para esta distribución. Para realizar el cálculo, denote el radio CEP por $r_{\mathrm{CEP}}$, plantee la ecuación $F(r_{\mathrm{CEP}})=0.5$ usando la FDA obtenida en el inciso anterior y despeje $r_{\mathrm{CEP}}$ en términos de $\sigma$.
    \item Las zonas de seguridad para evitar terrenos peligrosos (como cráteres o pendientes pronunciadas) suelen definirse como un radio de $3\sigma$ desde el centro. Calcule la probabilidad de que el rover aterrice fuera de esta zona segura.
    \item Supongamos que los sensores de a bordo detectan anomalías durante el descenso, y se sabe con certeza que el rover ha aterrizado a más de $1\sigma$ del centro (es decir, ya está fuera de la zona de error nominal). Calcule la probabilidad condicional de que el error radial supere los $2\sigma$.
    \item Preguntas conceptuales sobre el modelo:
    \begin{itemize}
        \item ¿Cuál es el significado físico y estadístico del parámetro $\sigma$ en el contexto del sistema de navegación del rover?
        \item ¿Por qué es matemática y físicamente apropiado utilizar la distribución de Rayleigh para modelar el error radial bajo las suposiciones dadas?
        \item ¿Qué representa operacionalmente el CEP para el equipo de ingeniería de la misión?
    \end{itemize}
\end{enumerate}


\section{Propagación de Error en el Sistema de Magnitudes Fotométricas}

En la astronomía observacional, la relación entre el flujo físico medido por un detector y la magnitud aparente reportada en los catálogos es altamente no lineal. Inspirados en los principios de fotometría y propagación de errores (\textit{Astronomical Measurement}, Capítulo sobre \textit{Photometry and Magnitudes}), analizaremos un problema complejo de transformación de variables aleatorias continuas.

La magnitud aparente $m$ de una fuente astrofísica se define en función de su flujo medido $f$ mediante la ecuación:
$$ m = -2.5 \log_{10}(f) + C $$
donde $C$ es el punto cero fotométrico (una constante instrumental). 

Supongamos que, debido a las características de ruido de un detector en el límite de sensibilidad o a la distribución estadística de una población de fuentes débiles, el flujo $f$ sigue una distribución exponencial con parámetro de tasa $\lambda > 0$. La función de densidad de probabilidad del flujo es:
$$ g(f) = \lambda e^{-\lambda f}, \quad \text{para } f > 0 $$

Para los cálculos puede utilizar la identidad
$$
\log_{10}(f)=\frac{\ln(f)}{\ln(10)}
$$
y definir la constante de escala $k=2.5/\ln(10)$. También puede emplear los siguientes resultados, donde $\gamma\approx 0.5772$ es la constante de Euler--Mascheroni:
$$
\int_{0}^{\infty}e^{-x}\,dx=1, \qquad
\int_{0}^{\infty}\ln(x)e^{-x}\,dx=-\gamma,
$$
$$
\int_{0}^{\infty}(\ln x)^2e^{-x}\,dx
=\gamma^2+\frac{\pi^2}{6}.
$$

\begin{enumerate}[a)]
    \item Reescriba la magnitud aparente en términos del logaritmo natural y de la constante $k$.
    \item Calcule el valor esperado $E(m)$. Plantee la integral correspondiente y realice el cambio de variable $x=\lambda f$, teniendo en cuenta que $\ln(f)=\ln(x)-\ln(\lambda)$.
    \item Calcule la varianza $V(m)$. Puede utilizar la definición $V(m)=E[(m-E(m))^2]$, el cambio de variable del inciso anterior y las identidades integrales proporcionadas.
    \item Sustituya $k=2.5/\ln(10)$ en la expresión obtenida para $V(m)$ y determine su valor numérico aproximado.
    \item Preguntas conceptuales sobre el modelo:
    \begin{itemize}
        \item ¿Cómo influyen el parámetro de tasa $\lambda$ y el punto cero instrumental $C$ en el valor esperado de la magnitud?
        \item ¿Depende la varianza de la magnitud de $\lambda$ o de $C$? Justifique su respuesta a partir del resultado obtenido.
        \item ¿Qué interpretación física tiene la transformación logarítmica del flujo en el sistema de magnitudes fotométricas?
    \end{itemize}
\end{enumerate}

\end{document}